\documentclass[11pt]{article}
\usepackage[a4paper,margin=27mm]{geometry}
\usepackage[T1]{fontenc}
\usepackage{lmodern}
\usepackage{amsmath,amssymb,amsthm}
\usepackage[hidelinks]{hyperref}
\hypersetup{pdftitle={Polynomiality of the Lecture Hall sequence via signed Dyck paths},pdfauthor={},pdfsubject={Lecture Hall polynomials and Stieltjes-Rogers moments},pdfcreator={}}
\newtheorem{theorem}{Theorem}[section]
\newtheorem{proposition}[theorem]{Proposition}
\newtheorem{lemma}[theorem]{Lemma}
\newtheorem{corollary}[theorem]{Corollary}
\theoremstyle{remark}
\newtheorem{remark}[theorem]{Remark}
\newcommand{\Z}{\mathbb Z}
\newcommand{\Q}{\mathbb Q}
\newcommand{\LH}{\ell}
\newcommand{\transpose}{\mathsf T}
\title{Polynomiality of the Lecture Hall sequence\\via signed Dyck paths}
\author{}
\date{October 13, 2026}

\begin{document}
\maketitle
\begin{abstract}
Katth\"an defined a sequence of rational functions by prescribing certain
Toeplitz minors and conjectured that every term is a polynomial. We identify
this sequence with a specialization of the Stieltjes--Rogers moment
polynomials: the down-step weight at height $h$ is $y_hy_{h+1}$, with a
negative sign precisely when $h\ge4$ is even. Classical Hankel determinant
identities, together with the signs arising from row reversal and the
uniqueness of the defining recurrence, establish this identification for
every index. Thus the polynomiality conjecture follows. We give a
self-contained derivation of the required classical identities and an
explicit coefficient formula indexed by compositions. The latter implies
that the $i$th term has exactly $2^{i-2}$ nonzero monomials for $i\ge2$.
We state the interpretation of two typographical inconsistencies in the
source and explain its support within that paper.
\end{abstract}

\section{The sequence and its moment representation}

Let $R=\Z[y_1,y_2,\ldots]$, where the variables are algebraically
independent, and let $K$ be its fraction field. For a sequence
$P=(P_1,P_2,\ldots)$ in $K$, define
\[
 M(P)_{r,c}=\begin{cases}-P_{c-r+1},&c\ge r,\\0,&c<r,\end{cases}
 \qquad r,c\ge1.
\]
We use the following square minors:
\begin{equation}\label{eq:E}
 E_i(P)=-\det M(P)_{
 \{1,\ldots,\lceil i/2\rceil\},
 \{\lfloor i/2\rfloor+1,\ldots,i\}},\qquad i\ge1.
\end{equation}
Set
\begin{equation}\label{eq:T}
 T_i=\prod_{j=1}^{i}y_j^{i+1-j}.
\end{equation}
The Lecture Hall sequence is the unique sequence
$\LH=(\LH_1,\LH_2,\ldots)$ in $K$ satisfying
\begin{equation}\label{eq:definition}
 E_i(\LH)=T_i\qquad(i\ge1).
\end{equation}
Existence and uniqueness are recalled in Lemma~\ref{lem:unique} below.
Katth\"an \cite[Definition 3.1 and Proposition 3.3]{Katthan} established
that these terms are integer Laurent polynomials and conjectured that
they are polynomials \cite[Conjecture 3.4]{Katthan}.

\subsection{Interpretation of the source definition}\label{sec:source}

Equation~\eqref{eq:E} uses the definition determined by the square-matrix
description and examples in \cite{Katthan}. On page 3 of the specified
arXiv version, both the displayed column set and the following sentence
start at $\lfloor i/2\rfloor$, without the $+1$. That set has one more
column than the prescribed number of rows; for $i=1$ it also includes
column 0. It therefore does not specify an ordinary determinant.

The intended square minor is fixed by three pieces of internal evidence.
First, the source describes the largest top-aligned square block, with
top-right entry $-P_i$, that contains none of the structural zeros below
the diagonal. If its size is $k$, its columns are $i-k+1,\ldots,i$, and
the condition is $i-k+1\ge k$. Its largest possible size is
$k=\lceil i/2\rceil$, giving exactly \eqref{eq:E}. Second, all four
determinants in Example 3.2 agree with \eqref{eq:E}. Third, Definition 4.1
defines $\LH_S$ using columns $S+1$ and identifies
$\LH_{\{\lfloor i/2\rfloor,\ldots,i-1\}}$ with $E_i(\LH)$.
This again gives the columns in \eqref{eq:E}.

There is a second inconsistency in Example 3.2 on page 4: the printed last
term of $\LH_4$ is $y_1y_2^2y_3^2y_4^2$. The fourth defining equation,
$\LH_2\LH_4-\LH_3^2=T_4$, instead gives
\begin{equation}\label{eq:ell4}
 \LH_4=y_1^2y_2^3(y_1+y_3)^2+y_1^2y_2^2y_3^2y_4.
\end{equation}
This is also the expression listed in Section 6, page 8, of the source.
We use the defining equations and this consistent expression. These are
explicit interpretations supported by the source itself; they are not
presented as a formally issued erratum.

\subsection{Statement and relation to existing theory}

A Dyck path of semilength $n$ starts at $(0,0)$, ends at $(2n,0)$, uses
steps $(1,1)$ and $(1,-1)$, and never goes below height 0. Give each up
step weight 1 and each down step from height $h$ weight $\alpha_h$. Write
$\mu_n(\alpha)$ for the sum of the products of step weights over all
such paths, with $\mu_0=1$. These are the Stieltjes--Rogers polynomials;
see \cite[Theorem 9.1(b)]{Sokal} and the path interpretation of continued
fractions in \cite{Flajolet}.

Define the signs and specialized weights by
\begin{equation}\label{eq:weights}
 \varepsilon_h=\begin{cases}-1,&h\ge4\text{ and }h\text{ is even},\\
 1,&\text{otherwise},\end{cases}
 \qquad \alpha_h=\varepsilon_h y_hy_{h+1}.
\end{equation}

\begin{theorem}\label{thm:main}
For the Lecture Hall sequence defined by \eqref{eq:E}--\eqref{eq:definition},
\begin{equation}\label{eq:main}
 \LH_{n+1}=y_1\mu_n(\alpha)\qquad(n\ge0),
\end{equation}
with the specialization \eqref{eq:weights}. In particular,
$\LH_i\in\Z[y_1,\ldots,y_i]$ for every $i\ge1$.
\end{theorem}

This proves Conjecture 3.4 of \cite{Katthan} under the source interpretation
in Section~\ref{sec:source}. The contribution presented here is the explicit
identification \eqref{eq:main} and its connection to the prescribed
Toeplitz minors. The underlying Hankel identities are classical: they
follow directly from \cite[Theorem 9.5(b,c)]{Sokal}. Section~\ref{sec:hankel}
includes a derivation to keep the argument self-contained. We do not
claim these general identities as new, nor assert priority for the
specific identification. The initial-algebra and SAGBI statements in
Conjecture 4.2 of \cite{Katthan} are outside the scope of this note.

\section{Classical Hankel identities for Dyck moments}\label{sec:hankel}

In this section the $\alpha_h$ are independent variables. All identities
are polynomial identities over $\Z$, so they remain valid under signed
specializations, including those with zero weights.

\begin{lemma}[Hankel products]\label{lem:hankel}
For every $m\ge1$, the Dyck moments satisfy
\begin{align}
 H_m^{(0)}&:=\det[\mu_{i+j}]_{0\le i,j<m}
   =\prod_{h=1}^{m-1}(\alpha_{2h-1}\alpha_{2h})^{m-h},\label{eq:H0}\\
 H_m^{(1)}&:=\det[\mu_{i+j+1}]_{0\le i,j<m}
   =\alpha_1^m\prod_{h=1}^{m-1}
       (\alpha_{2h}\alpha_{2h+1})^{m-h}.\label{eq:H1}
\end{align}
Empty products are 1.
\end{lemma}

\begin{proof}
Group the steps of a Dyck path into consecutive pairs and divide its
heights at even times by 2. The resulting Motzkin path has up-step weight
1 and down-step weight
\[
 \lambda_h=\alpha_{2h-1}\alpha_{2h}\qquad(h\ge1).
\]
The two possible horizontal steps at height $h\ge1$ have weights
$\alpha_{2h}$ and $\alpha_{2h+1}$; at height 0 only the latter is allowed.
Their total weights are therefore
\[
 b_0=\alpha_1,\qquad b_h=\alpha_{2h}+\alpha_{2h+1}\quad(h\ge1).
\]
Let $J$ be the infinite matrix indexed by the nonnegative integers with
$J_{h,h}=b_h$ and $J_{h,h+1}=1$ for $h\ge0$, and
$J_{h,h-1}=\lambda_h$ for $h\ge1$; all other entries are zero.
Matrix multiplication enumerates these paths,
so $\mu_n=(J^n)_{0,0}$. Each relevant entry of a fixed power involves only
finitely many indices.

Set $d_0=1$, $d_h=\lambda_1\cdots\lambda_h$, and
$D=\operatorname{diag}(d_0,d_1,\ldots)$. Directly,
$JD=DJ^\transpose$, and hence
\begin{equation}\label{eq:balance}
 (J^j)_{h,0}=d_h(J^j)_{0,h}\qquad(h,j\ge0).
\end{equation}
No division by $d_h$ is used. Fix $m\ge1$ and put
$L_{i,h}=(J^i)_{0,h}$ for $0\le i,h<m$. Since a path of length $i$
cannot reach a height greater than $i$, and the unique path reaching
height $i$ has weight 1, $L$ is unit lower triangular.

Let $D_m$ and $J_m$ denote the leading $m\times m$ blocks of $D$ and $J$.
Using \eqref{eq:balance} gives
\begin{equation}\label{eq:factor0}
 [\mu_{i+j}]_{0\le i,j<m}=LD_mL^\transpose.
\end{equation}
Indeed, the intermediate height in $J^iJ^j$ is at most $\min(i,j)$.
For the shifted matrix, write
\begin{align*}
 \mu_{i+j+1}
 &=\sum_{a,b\ge0}(J^i)_{0,a}J_{a,b}(J^j)_{b,0}\\
 &=\sum_{a,b=0}^{m-1}L_{i,a}(J_m)_{a,b}d_bL_{j,b}.
\end{align*}
The truncation is exact: a nonzero summand requires $a\le i<m$ and
$b\le j<m$. Consequently,
\begin{equation}\label{eq:factor1}
 [\mu_{i+j+1}]_{0\le i,j<m}=LJ_mD_mL^\transpose.
\end{equation}
Taking determinants in \eqref{eq:factor0} proves \eqref{eq:H0}.

For \eqref{eq:H1}, let $\delta_m=\det J_m$ and $\delta_0=1$.
Expansion of a tridiagonal determinant yields
\[
 \delta_1=\alpha_1,\qquad
 \delta_m=b_{m-1}\delta_{m-1}
             -\lambda_{m-1}\delta_{m-2}\quad(m\ge2).
\]
Induction gives $\delta_m=\prod_{h=1}^{m}\alpha_{2h-1}$: after using the
induction hypothesis, the summand containing $\alpha_{2m-2}$ from
$b_{m-1}$ cancels the second term of the recurrence. This is a polynomial
cancellation and does not require any factor to be invertible.
Taking determinants in \eqref{eq:factor1} now gives
\[
 H_m^{(1)}=\left(\prod_{h=1}^m\alpha_{2h-1}\right)
              \prod_{h=1}^{m-1}(\alpha_{2h-1}\alpha_{2h})^{m-h},
\]
which rearranges to \eqref{eq:H1}. For $m=1$, the two formulas read
$H_1^{(0)}=1$ and $H_1^{(1)}=\alpha_1$, as required.
\end{proof}

For comparison, the two Stieltjes--Rogers factorizations in
\cite[Theorem 9.5(b,c)]{Sokal} have diagonal entries
$\prod_{j=1}^{2h}\alpha_j$ and $\prod_{j=1}^{2h+1}\alpha_j$,
respectively, for $h\ge0$. Their leading principal determinants give
exactly \eqref{eq:H0} and \eqref{eq:H1}. Positivity assumptions from
analytic moment problems are not needed here.

\section{Identification with the Lecture Hall sequence}

\begin{lemma}[Triangular recurrence]\label{lem:unique}
There is exactly one sequence in $K$ satisfying \eqref{eq:definition}.
\end{lemma}
\begin{proof}
The first two equations give $\LH_1=y_1$ and $\LH_2=y_1^2y_2$.
For $i\ge3$, put $k=\lceil i/2\rceil$. The entry $-P_i$ occurs just
once in the minor in \eqref{eq:E}, in its first row and last column.
Its cofactor, after shifting both remaining index sets down by 1, is
the minor defining $E_{i-2}$. Keeping track of the determinant signs
gives
\[
 E_i(P)=(-1)^kE_{i-2}(P)P_i+Q_i(P_1,\ldots,P_{i-1})
\]
for a polynomial $Q_i$. After the previous equations have been solved,
the coefficient of $P_i$ is the nonzero monomial $(-1)^kT_{i-2}$.
It is invertible in $K$, so each equation has exactly one solution for
its new term. Induction proves both existence and uniqueness.
\end{proof}

\begin{proof}[Proof of Theorem~\ref{thm:main}]
Let $P_{n+1}=y_1\mu_n(\alpha)$ with \eqref{eq:weights}. Define
\[
 s_m=(-1)^{1+m(m+1)/2}.
\]
Reverse the row order in the square minor \eqref{eq:E} and factor a
minus sign from each row. The row reversal has sign
$(-1)^{m(m-1)/2}$, and the negative sign outside the determinant must
also be included. Thus for every $m\ge1$,
\begin{align}
 E_{2m-1}(P)&=s_m\det[P_{r+s+1}]_{0\le r,s<m}
             =s_my_1^m H_m^{(0)},\label{eq:odd}\\
 E_{2m}(P)&=s_m\det[P_{r+s+2}]_{0\le r,s<m}
             =s_my_1^m H_m^{(1)}.\label{eq:even}
\end{align}

Apply Lemma~\ref{lem:hankel}. In \eqref{eq:H0}, both
$\alpha_{2h-1}$ and $\alpha_{2h}$ have exponent $m-h$. After the
substitution \eqref{eq:weights} and multiplication by $y_1^m$, the
exponent of $y_j$ is $2m-j$ for $1\le j\le2m-1$. In \eqref{eq:H1},
$\alpha_1$ has exponent $m$, while $\alpha_{2h}$ and
$\alpha_{2h+1}$ have exponent $m-h$. The resulting exponent of $y_j$
is $2m+1-j$ for $1\le j\le2m$. These are precisely the exponents
in $T_{2m-1}$ and $T_{2m}$.

For $m\ge2$, the sign in each product is
\[
 (-1)^{\sum_{h=2}^{m-1}(m-h)}
   =(-1)^{(m-2)(m-1)/2}=s_m,
\]
since the last exponent and $1+m(m+1)/2$ differ by $2m$.
For $m=1$ the two products give $y_1$ and $y_1^2y_2$ directly, and
$s_1=1$. Hence
\[
 y_1^mH_m^{(0)}=s_mT_{2m-1},\qquad
 y_1^mH_m^{(1)}=s_mT_{2m}.
\]
Equations~\eqref{eq:odd}--\eqref{eq:even} therefore give
$E_i(P)=T_i$ for every $i\ge1$. Lemma~\ref{lem:unique} implies
$P_i=\LH_i$. Finally, a Dyck path of semilength $n$ has height at
most $n$, so $y_1\mu_n(\alpha)$ is a finite sum in
$\Z[y_1,\ldots,y_{n+1}]$.
\end{proof}

\begin{remark}
The uniqueness argument takes place in the fraction field of independent
variables. This stepwise division argument cannot be applied directly
to specializations at which a divisor vanishes.
This does not restrict the polynomial identity in Theorem~\ref{thm:main},
which can itself be evaluated at arbitrary values.
\end{remark}

\section{Coefficients and generating function}

\begin{proposition}\label{prop:coefficients}
For every $n\ge1$,
\begin{equation}\label{eq:composition}
 \LH_{n+1}=y_1
 \sum_{\substack{r\ge1,\ m_1+\cdots+m_r=n\\m_1,\ldots,m_r\ge1}}
 \left(\prod_{j=1}^{r-1}
       \binom{m_j+m_{j+1}-1}{m_{j+1}}\right)
 (-1)^{m_4+m_6+\cdots}
 \prod_{j=1}^{r}(y_jy_{j+1})^{m_j}.
\end{equation}
An absent $m_j$ is read as zero. Consequently $\LH_i$ has exactly
$2^{i-2}$ nonzero monomials for every $i\ge2$.
\end{proposition}
\begin{proof}
For a nonempty Dyck path, let $m_j$ be the number of down steps from
height $j$, and let $r$ be its maximum height. Every $m_j$ with
$1\le j\le r$ is positive and their sum is $n$.

Use the usual contour correspondence between Dyck paths and plane
rooted trees. Then $m_j$ is the number of vertices at depth $j$.
Given the vertices at depth $j$ in their left-to-right order, specifying
the number of children of each is a weak composition of $m_{j+1}$ into
$m_j$ parts. There are
$\binom{m_j+m_{j+1}-1}{m_{j+1}}$ choices. These choices can be made
independently at each depth and determine the ordered tree uniquely.
The root has $m_1$ children, contributing a factor of 1. Thus the product
of binomial coefficients in \eqref{eq:composition} counts paths with
the prescribed profile. Their common weight is
$\prod_j\alpha_j^{m_j}$, which proves the formula by
Theorem~\ref{thm:main}.

Different compositions produce different monomials. If $e_j$ is the
exponent of $y_j$, then $m_1=e_1-1$ and
$m_j=e_j-m_{j-1}$ for $j\ge2$, with missing parts set to zero.
Thus the profile can be recovered from the monomial. Each coefficient
has strictly positive absolute value, so no profile disappears by
cancellation. There are $2^{n-1}$ compositions of $n$, giving the claim
with $i=n+1$.
\end{proof}

For example, the first terms are $\LH_1=y_1$, $\LH_2=y_1^2y_2$,
$\LH_3=y_1^2y_2^2(y_1+y_3)$, the expression \eqref{eq:ell4}, and
\begin{align*}
 \LH_5={}&y_1^2y_2^4(y_1+y_3)^3
       +2y_1^2y_2^3y_3^2y_4(y_1+y_3)\\
       &+y_1^2y_2^2y_3^2y_4^2(y_3-y_5).
\end{align*}
In particular, polynomiality does not mean coefficientwise
nonnegativity.

The path interpretation also gives the formal Stieltjes continued
fraction
\begin{equation}\label{eq:fraction}
 \sum_{n\ge0}\mu_n(\alpha)z^n
 =\cfrac{1}{1-\cfrac{\alpha_1z}{1-\cfrac{\alpha_2z}{1-\cfrac{\alpha_3z}{\ddots}}}}.
\end{equation}
Indeed, if $F_h$ enumerates excursions above height $h$, the first-return
decomposition gives $F_h=1+\alpha_{h+1}zF_{h+1}F_h$.
The denominator in $F_h=(1-\alpha_{h+1}zF_{h+1})^{-1}$ has constant
term 1. Moreover, $[z^n]F_0$ only uses heights up to $n$, so the infinite
fraction is well-defined coefficient by coefficient. There is no
analytic convergence or sign condition.

\section{Exact finite checks and scope}

The accompanying \texttt{verify.py} provides three exact checks. It
generates the path polynomials by stepwise dynamic programming and
substitutes them into the original negative Toeplitz minors
\eqref{eq:E}, checking all symbolic identities for $1\le i\le10$.
It also checks the defining equations through $i=40$ for each of the
three integer assignments
\[
 y_j=1,\qquad y_j=j,\qquad
 y_j=(-1)^j\bigl((j\bmod5)+1\bigr).
\]
Finally, it compares the path polynomials with the composition formula
through $\mu_{12}$, equivalently $\LH_{13}$. The recorded output is
\texttt{verification.json}; the calculation uses exact integer and
polynomial arithmetic. The last comparison does not claim a symbolic
Toeplitz determinant check through index 13.

These finite checks aid reproducibility and detect indexing or sign
errors. The all-index statement follows from the preceding proof, not
from extrapolation of the checks. The argument establishes polynomiality
and the displayed coefficient formula; it does not establish the SAGBI
or initial-algebra conjectures of \cite{Katthan}.

\paragraph{AI assistance.}
This work was developed with AI assistance. The mathematical argument
is given in full, and exact finite verification code is supplied for
reproducibility.

\begin{thebibliography}{9}

\bibitem{Katthan}
Lukas Katth\"an,
\emph{The Lecture Hall Cone as a toric deformation},
arXiv:1809.01377v1 (2018).
\url{https://arxiv.org/abs/1809.01377v1}.

\bibitem{Sokal}
Alan D. Sokal,
\emph{A simple algorithm for expanding a power series as a continued fraction},
Expositiones Mathematicae \textbf{41} (2023), 245--287.
\href{https://doi.org/10.1016/j.exmath.2022.12.001}{doi:10.1016/j.exmath.2022.12.001}.
Also arXiv:2206.15434.

\bibitem{Flajolet}
Philippe Flajolet,
\emph{Combinatorial aspects of continued fractions},
Discrete Mathematics \textbf{32} (1980), 125--161.
\url{https://algo.inria.fr/flajolet/Publications/Flajolet80b.pdf}.

\end{thebibliography}
\end{document}
